考虑一无源、线性、互易的对称双端口光学微腔，中心角频率为 \(\omega_0>0\)。微腔仅支持一个共振模，其慢变振幅 \(a(t)\) 归一化为储能 \(|a|^2\)，两端口入射、出射行波 \(s_{1+},s_{2+}\) 和 \(s_{1-},s_{2-}\) 归一化为功率 \(|s|^2\)。内部吸收率为 \(\gamma_i>0\)，向两端口的总辐射振幅衰减率为 \(\gamma_e>0\)。定义对称和反对称通道 \(s_{\pm}^{(s)}=(s_{1\pm}+s_{2\pm})/\sqrt2\)、\(s_{\pm}^{(a)}=(s_{1\pm}-s_{2\pm})/\sqrt2\)。在背景直反射相位为 \(-1\) 的规范下，时域耦合模方程为
\[
 \dot a=(-i\omega_0-\gamma_i-\gamma_e)a+\sqrt{2\gamma_e}\,s_+^{(s)},\quad
 s_-^{(s)}=-s_+^{(s)}+\sqrt{2\gamma_e}\,a,\quad
 s_-^{(a)}=-s_+^{(a)}.
\]
所有稳态行波采用 \(e^{-i\omega t}\) 约定，记 \(\Delta=\omega-\omega_0\)。
(1) 由时域方程求双端口散射矩阵 \(S(\omega)\)，对角化其对称和反对称本征通道；验证 \(\gamma_i=0\) 时散射幺正且极点位于下半复频平面。
(2) 推导相干完美吸收（CPA）所需的频率、损耗匹配及入射波相位条件；证明有非零反对称入射时不可能实现全吸收。
(3) 求对称通道的吸收率线型、半高全宽及加载品质因子；从腔内储能与内耗功率独立验证能量守恒和吸收率表达式。
(4) 两端口输入相同功率且相位差为 \(\phi\) 时，求总吸收率随 \(\phi\) 的变化、最优与最差相位；再推导单端入射的最大吸收率，解释为什么双端口相干驱动可以达到而单端驱动不能达到100%吸收。
